\section{From spectral positivity to a quadratic reserve}
\label{sec:preprocessing}

Recall the weighted scalar expression $\Dcal$, and define its associated
weighted divergence by
\begin{equation}\label{pre:drift-scalar}
 \Dcal(h,t)=\Scal_h+2\Delta_h t-|\dd t|_h^2,
 \qquad
 \diver_{h,t}Y=\diver_hY-\dd t(Y).
\end{equation}
The spectral hypothesis first supplies a function $w$ with
$\Dcal(g,w)\ge1$. The difficulty is that it supplies no uniform control
of $w$. Our aim in this section is to enlarge the metric and replace $w$
by a new function for which the scalar inequality contains a positive
quadratic term. This term will pay for the later distance enlargements.

The passage from a quadratic-form inequality to a positive solution is
part of Allegretto--Piepenbrink theory
\cite{Allegretto1974,MossPiepenbrink1978}. We include the exhaustion
argument in the form needed here.

\begin{lemma}[Positive supersolution]\label{pre:supersolution}
Let $(M,g)$ be a connected complete smooth manifold without boundary
satisfying the normalized spectral bound in Theorem~\ref{thm:main}.
There is a smooth positive function $v$ on $M$ such that
\[
 -4\Delta_gv+\Scal_gv\ge v.
\]
Consequently $w=-2\log v$ is smooth and satisfies $\Dcal(g,w)\ge1$.
\end{lemma}

\begin{proof}
If $M$ is compact, take a positive first eigenfunction of
$-4\Delta_g+\Scal_g$. For noncompact $M$, choose a nested exhaustion
by relatively compact connected smooth domains $\Omega_j$, with
$\overline\Omega_j\subset\Omega_{j+1}$, all containing a fixed point $o$.
Let $v_j$ be a positive Dirichlet first eigenfunction on $\Omega_j$,
normalized by $v_j(o)=1$, and let $\lambda_j$ be its eigenvalue.
These eigenfunctions exist by the variational principle, compact Sobolev
embedding on each domain, elliptic regularity, and the strong maximum
principle. The potential is smooth and bounded on each domain closure.
Approximation in the Dirichlet form domain shows that $\lambda_j\ge1$.
One fixed nonzero test function supported in $\Omega_1$ bounds all
$\lambda_j$ from above.

After passing to a subsequence, $\lambda_j\to\lambda\ge1$.
On every fixed compact subset of $M$, interior Harnack inequalities
and finite chains of coordinate balls bound $v_j$ above and below,
using the normalization at $o$. Interior elliptic estimates and a
diagonal subsequence then give a smooth positive limit $v$ satisfying
$-4\Delta_gv+\Scal_gv=\lambda v$. All estimates use only fixed compact
subsets; no uniform geometry at infinity is required. Finally,
\[
 2\Delta_g(-2\log v)-|\dd(-2\log v)|_g^2
       =-4\frac{\Delta_gv}{v},
\]
which proves the assertion for $w$.
\end{proof}

The reserve construction applies to arbitrary smooth functions, not
only to the function supplied by the spectral hypothesis.

\begin{proposition}[Quadratic reserve]\label{pre:reserve}
For every integer $n\ge4$ there is a constant $C(n)<\infty$ with the
following property. If $h$ is a complete smooth metric on a connected
boundaryless $n$-manifold $M$ and $w\in C^\infty(M)$, there are a smooth
complete metric $h^+\ge h$ and a smooth real-valued function $t^+$ such that
\begin{equation}\label{pre:reserve-bound}
 \Dcal(h^+,t^+)\ge \Dcal(h,w)+(t^+)^2-C(n)
 \quad\hbox{on }M.
\end{equation}
One may take $C(n)=100n+203$.
\end{proposition}

We shall use the metric induced by a graph of height $u$ in
$(M,h)\times\R$, prescribing its weighted mean curvature by
\[
 \diver_{h,w}\left(\frac{\nabla_hu}{\sqrt{1+|\dd u|_h^2}}\right)
     =p,\qquad p=10(u-w).
\]
The pressure measures the discrepancy between the graph height and the
original function, even when both are unbounded.
For a candidate height $u$, put $s=|\dd u|_h$ and
$l=(1+s^2)^{-1/2}$. The induced metric is $h_u=h+\dd u\otimes\dd u$,
and its volume density is $l^{-1}\dd\vol_h$. Thus $t_u=t-\log l$ is
exactly the change preserving the weighted density:
\[
 e^{-t_u}\dd\vol_{h_u}=e^{-t}\dd\vol_h.
\]
For $t=w$ and pressure $p=10(u-w)$, our final drift will satisfy
\[
 t^+=t_u-u=-p/10-\log l.
\]
This is the quantity whose square the graph equation must supply.
The scalar calculation below produces a positive $p^2$ term, together
with slope and derivative terms that we will also estimate. We first
construct a smooth solution of the graph equation; the scalar comparison
and the square completion will then be carried out consecutively.

\subsection{The angle equation and elliptic estimates}
\label{pre:elliptic-tools}

To construct the prescribed graph, we need an a priori gradient estimate
on a possibly noncompact manifold. The useful
gradient variable is the vertical component of the graph normal: a
positive lower bound for this component is an upper bound for the slope.
We record the equation it satisfies, allowing angle-dependent pressure
for use in the later local construction. The use of normal variation
and comparison to estimate graph gradients has a classical predecessor
in Korevaar's treatment of prescribed mean curvature~\cite{Korevaar1984};
we give the weighted calculation and all estimates used here.

\begin{lemma}[Graph angle equation]\label{pre:angle}
Let $u$ be smooth on an open subset of a Riemannian manifold $(M,h)$.
On its graph $\Gamma$ in $(M,h)\times\R$, with vertical coordinate $z$,
put
\[
 \omega=(1+|\dd u|_h^2)^{-1/2},\qquad
 \nu=(\omega\nabla_hu,-\omega),\qquad Z=-\partial_z.
\]
Let $t$ be a smooth function on the base, extended constantly in $z$,
and let $X=-\nabla t$. Suppose that a smooth function $p(x,z,\omega)$,
defined near the graph data, satisfies
\[
 H_\Gamma+\langle X,\nu\rangle=p(x,u(x),\omega).
\]
Here $H_\Gamma$ is the trace of the second fundamental form
$\mathrm{II}=\nabla\nu$. Then
\begin{equation}\label{pre:angle-equation}
 \begin{split}
 \mathcal L\omega&=-p_z,\\
 \mathcal L={}&\Delta_\Gamma
  +\langle X^\top-p_\omega Z^\top,\nabla_\Gamma\,\cdot\,\rangle\\
 &+|\mathrm{II}|^2+\Ric_{\mathrm{amb}}(\nu,\nu)
       +\Hess_{\mathrm{amb}}t(\nu,\nu)
       +\partial_\nu^{\mathrm{fix}\,\omega}p.
 \end{split}
\end{equation}
The superscript $\top$ denotes tangential projection to $\Gamma$;
the final derivative differentiates the ambient variables while
holding the angle argument of $p$ fixed.
\end{lemma}

\begin{proof}
Extending the displayed normal independently of $z$ gives
$H_\Gamma=\diver_h(\omega\nabla_hu)$. For a normal variation with speed
$V$, the variation formulas in our convention are
\[
 \dot\nu=-\nabla_\Gamma V,\qquad
 \dot H_\Gamma=-\Delta_\Gamma V
       -\bigl(|\mathrm{II}|^2+\Ric_{\mathrm{amb}}(\nu,\nu)\bigr)V.
\]
Since $X=-\nabla t$ and $\omega=\langle Z,\nu\rangle$, one has
\[
 \begin{split}
 \frac{\dd}{\dd a}\Big|_{a=0}\langle X,\nu\rangle
   &=-\Hess_{\mathrm{amb}}t(\nu,\nu)V
                          -\langle X^\top,\nabla_\Gamma V\rangle,\\
 \dot p&=\partial_\nu^{\mathrm{fix}\,\omega}p\,V
                       -p_\omega\langle Z^\top,\nabla_\Gamma V\rangle.
 \end{split}
\]
Thus the linearization of $H_\Gamma+\langle X,\nu\rangle-p$ is
$-\mathcal LV$. Translation in the direction $Z=-\partial_z$ has normal
speed $\omega$. It leaves the geometric left-hand side and the angle
unchanged, while its derivative of the residual is $+p_z$.
The tangential component differentiates a residual which vanishes
identically on $\Gamma$. Hence $-\mathcal L\omega=p_z$.
\end{proof}

We will use the following local elliptic facts in both graph constructions.
In coordinates, the equations have the form
\begin{equation}\label{pre:quasilinear-equation}
 \partial_i\mathfrak A^i(x,Du)=\mathfrak B(x,u,Du),
\end{equation}
where the volume density may be included in $\mathfrak A$ and
$\mathfrak B$. Suppose these functions are smooth on a neighborhood
of a compact range of $(x,u,Du)$, with bounded derivatives there, and
that $\mathfrak A^i_{p_m}$ is uniformly elliptic on that range.
For a family of equations, the stated coefficient and ellipticity
bounds are assumed uniform over the parameters under consideration.
Uniform height and gradient bounds for smooth solutions then give
interior bounds of every order. For completeness, differentiating in
$x_j$ and writing $b=\partial_j u$ gives
\begin{equation}\label{pre:differentiated-equation}
 \partial_i\bigl(\mathfrak A^i_{p_m}\partial_m b\bigr)
       -\mathfrak B_{p_m}\partial_m b
   =\mathfrak B_u b+\mathfrak B_{x_j}
                       -\partial_i\mathfrak A^i_{x_j}.
\end{equation}
Subscripts $x_j$ on the coefficient functions mean explicit derivatives,
with the other arguments held fixed. Thus this is a uniformly elliptic
divergence equation with bounded coefficients, bounded lower-order
terms, bounded scalar data, and the divergence of bounded vector data.
The interior De Giorgi--Nash estimate~\cite[Theorem~8.24]{GilbargTrudinger2001}
gives a uniform H\"older bound
for $b$, and hence a $C^{1,\alpha}$ bound for $u$, for some $\alpha>0$.
The nondivergence form of \eqref{pre:quasilinear-equation} now has
H\"older coefficients and right-hand side. Interior Schauder estimates
\cite[Theorem~6.2 and Corollary~6.3]{GilbargTrudinger2001}
give $C^{2,\alpha}$ bounds, and differentiation yields bounds of all
higher orders; compare also~\cite[Theorem~13.1]{GilbargTrudinger2001}
for this differentiation argument. On a closed manifold a finite coordinate cover makes
these estimates global.

In particular, a continuation on a fixed closed manifold is closed
once uniform height and gradient bounds keep the solutions in a compact
subset of the open equation domain. It is open when the linearization
has strictly negative zeroth-order coefficient. Indeed, for the resulting
smooth uniformly elliptic operator, interpolation with $\Delta-1$,
the maximum principle, and Schauder estimates give invertibility
$C^{2,\alpha}\to C^{0,\alpha}$. The starting operator $\Delta-1$ is
invertible by the energy method and elliptic regularity. The implicit
function theorem then gives openness, and elliptic regularity gives
smooth solutions. All constants in this discussion may depend on the
individual compact manifold or the fixed coordinate charts.

\subsection{A global graph with linear discrepancy pressure}

We now solve the graph equation without imposing any global size
condition on $w$. The positive derivative of the pressure with respect to height
will provide both comparison and the angle bound.

\begin{lemma}[Global prescribed graph]\label{pre:global-graph}
Let $h$ be a complete smooth metric on a connected manifold $M$ without
boundary, and let $w\in C^\infty(M)$. With $c=10$, there is a smooth
function $u$ on $M$ satisfying
\begin{equation}\label{pre:global-equation}
 \diver_{h,w}\left(\frac{\nabla_hu}{\sqrt{1+|\dd u|_h^2}}\right)
       =c(u-w).
\end{equation}
\end{lemma}

\begin{proof}
First suppose $M$ is closed. Replace $w$ by $\tau w$ in both occurrences
of \eqref{pre:global-equation}, with $0\le\tau\le1$. The zero function
solves the equation at $\tau=0$. Comparison with constants gives
$|u|\le\max_M|w|$ along the continuation. In ordinary mean-curvature
form, the equation has $t=\tau w$ and $p=c(z-\tau w(x))$ in
Lemma~\ref{pre:angle}; in particular $p_z=c$ and $p_\omega=0$.
The zeroth-order coefficient of $\mathcal L$, after discarding
$|\mathrm{II}|^2$, is bounded below by $-C$, uniformly in $\tau$ and
the solution. This uses only the compact bounds for the ambient Ricci
tensor, $\Hess w$, and $\dd w$, since
\[
 \partial_\nu^{\mathrm{fix}\,\omega}p
      =c\bigl(\nu_z-\tau\dd w(\nu_h)\bigr).
\]
At a minimum of $\omega$, Equation~\eqref{pre:angle-equation} gives
\[
 -c=\mathcal L\omega\ge-C\omega.
\]
Thus $\omega\ge c/C>0$, after increasing $C$ if necessary, and the
gradient is uniformly bounded. The mean-curvature flux has principal
eigenvalues $(1+s^2)^{-1/2}$ orthogonal to $\nabla u$ and
$(1+s^2)^{-3/2}$ in its direction, where $s=|\dd u|_h$.
The gradient bound therefore gives uniform ellipticity. The
linearization has zeroth-order coefficient $-c$; the preceding elliptic
facts give existence at $\tau=1$.

For noncompact $M$, choose closed manifolds extending successively larger
relatively compact regions, with extensions of $h$ and $w$ agreeing
with the original data near each region closure. To obtain them, take
smooth compact domain enlargements, modify the data outside the region
to be product data in a boundary collar, and double. This construction
does not require orientability. Solve the equation on each closed
extension. It remains to obtain estimates on every fixed compact
subset of the original manifold that do not depend on the extension.

Fix a sufficiently small normal coordinate ball $B$ of radius $R$,
whose closure is contained in a normal coordinate neighborhood. We
consider only extensions agreeing with the original data near
$\overline B$. Write $r$ for radial distance from its center.
Choose a nonnegative smooth radial function $b$ on $B$, constant near
the center and equal to $(R-r)^{-1}$ near the boundary. The expression
$\diver_{h,w}(\nabla_h b/\sqrt{1+|\dd b|_h^2})$ is bounded on $B$.
In fact, near the boundary,
\[
 \frac{b'}{\sqrt{1+(b')^2}}=(1+(R-r)^4)^{-1/2},\qquad
 \left(\frac{b'}{\sqrt{1+(b')^2}}\right)'
       =\frac{2(R-r)^3}{(1+(R-r)^4)^{3/2}},
\]
and the remaining terms use bounded $\Delta_hr$ and $\dd w$ there.
Taking $C$ sufficiently large makes $C+b$ an upper solution and
$-C-b$ a lower solution of \eqref{pre:global-equation} on $B$.
Both barriers diverge with the appropriate sign at the boundary.
Each solution under consideration is smooth across $\overline B$,
so a violation attains an interior extremum. At such an extremum the
gradients agree, ellipticity compares the second-order terms, and the
strict increase of $c(u-w)$ in $u$ gives a contradiction. We obtain
uniform height bounds on every smaller ball.

To bound the gradient there, choose a nonnegative smooth cutoff
$\zeta_0$ compactly supported in $B$ and positive on that smaller ball.
Use Lemma~\ref{pre:angle} with $t=w$ and $p=c(z-w(x))$.
The height bound on $\supp\zeta_0$ bounds $p$ and hence $H_\Gamma$.
Viewing $\zeta_0$ as a function on the graph, one has
\[
 \Delta_\Gamma\zeta_0
   =\tr_{T\Gamma}\Hess_{\mathrm{amb}}\zeta_0
       -H_\Gamma\langle\nabla\zeta_0,\nu\rangle.
\]
Thus its intrinsic Laplacian and gradient are bounded independently of
the slope. The drift $X^\top$ is bounded, and the zeroth-order
coefficient of $\mathcal L$ is bounded below after discarding
$|\mathrm{II}|^2$.

Let $A=\max(\zeta_0/\omega)>0$. At a point attaining this maximum,
$\omega-\zeta_0/A$ has value and gradient zero and nonnegative
Laplacian. Therefore
\begin{equation}\label{pre:cutoff-angle}
 -c=\mathcal L\omega
       \ge A^{-1}\mathcal L\zeta_0\ge-C'/A.
\end{equation}
The last bound is uniform on the fixed ball: the possibly unbounded
term $|\mathrm{II}|^2\zeta_0$ is nonnegative, and every other term has
the bounds just established. It follows that $A\le C'/c$. Since
$\zeta_0$ has a positive lower bound on any compact subball on which
it is positive, this is the required interior gradient bound.

The elliptic estimates above now give uniform bounds of every order
on fixed compact subsets. A diagonal subsequence of the extension
solutions converges smoothly locally on $M$ and solves
\eqref{pre:global-equation}. None of the constants in this limiting
argument requires a uniform geometric bound at infinity.
\end{proof}

\subsection{Scalar comparison for a graph enlargement}
\label{pre:scalar-subsection}

The global graph is now available, and we turn to its scalar curvature.
For the later localized construction we allow a more general metric
which stretches only the line normal to a level set of the height function.
Writing this length multiplier as $1/l$, we again use $t_u=t-\log l$
to preserve the weighted density. The following calculation extends the
graph comparison underlying the local deformation in
\cite[Section~3.5, Lemma~3.4]{OpenAIWidth}.

\begin{lemma}[Scalar graph comparison]\label{pre:graph-comparison}
Let $h$ be a smooth metric on an $n$-manifold, let $t,u$ be smooth
functions, and work on an open set where $s=|\dd u|_h>0$.
Let $\theta=\theta(s)$ be a smooth function with $0<\theta<1$, and put
\[
 l=\sqrt{1-\theta},\qquad \beta=\log l,
 \qquad k_1=\sqrt\theta,\qquad e=\nabla_hu/s.
\]
Assume $0\le-s\beta_s\le2$. Define
\begin{equation}\label{pre:graph-update}
 h_u=h+\frac{\theta}{(1-\theta)s^2}\,\dd u\otimes\dd u,
 \qquad t_u=t-\beta.
\end{equation}
Let $J$ be the second fundamental form of the level sets of $u$ with
normal $e$, let $H=\tr J$, and set
\[
 S=H-e(t),\qquad K=\frac{l^2}{\theta}e(\beta),
 \qquad M_1=S-K.
\]
Then $K=-e(\log k_1)$, and
\begin{equation}\label{pre:weighted-divergence}
 \diver_{h,t}(k_1e)=k_1M_1.
\end{equation}
Moreover,
\begin{equation}\label{pre:scalar-comparison}
 \begin{split}
 \Dcal(h_u,t_u)-\Dcal(h,t)&\ge\theta E,\\
 E&=|J|^2+M_1^2+K^2/l^2-2M_1K+2e(M_1).
 \end{split}
\end{equation}
All norms of level tensors in this formula use the metric induced by $h$.
\end{lemma}

\begin{proof}
Use $q=u$ as a transverse coordinate and transport the remaining
coordinates along the flow orthogonal to the levels. Locally,
\[
 h=v_1^2\,\dd q^2+\gamma_q,\qquad
 v_1=s^{-1},\qquad e=v_1^{-1}\partial_q,
 \qquad J=\frac{1}{2v_1}\partial_q\gamma_q.
\]
The Gauss equation and the mean-curvature variation formula give
\[
 \Scal_h=\Scal_{\gamma_q}-|J|^2-H^2-2e(H)
                         -2\frac{\Delta_{\gamma_q}v_1}{v_1}.
\]
Indeed, the normal variation with speed $v_1$ has mean-curvature
derivative $-\Delta_{\gamma_q}v_1-(|J|^2+\Ric_h(e,e))v_1$.
The Laplacian of $t$ is
\[
 \Delta_ht=\Delta_{\gamma_q}t+e(e(t))+He(t)
       +\langle\dd_\gamma\log v_1,\dd_\gamma t\rangle.
\]
Consequently
\begin{equation}\label{pre:foliation-scalar}
 \begin{split}
 \Dcal(h,t)={}&-|J|^2-S^2-2e(S)
       +\Scal_{\gamma_q}+2\Delta_{\gamma_q}t-|\dd_\gamma t|^2\\
 &-2\frac{\Delta_{\gamma_q}v_1}{v_1}
       +2\langle\dd_\gamma\log v_1,\dd_\gamma t\rangle.
 \end{split}
\end{equation}
Here $\dd_\gamma$ denotes tangential differentiation on a level set.

For the new metric, $v_1$ is replaced by $v_1/l$ and $t$ by $t-\beta$.
The resulting change in the last four tangential terms of
\eqref{pre:foliation-scalar} is
\[
 2\langle\dd_\gamma\log v_1,\dd_\gamma\beta\rangle
       -|\dd_\gamma\beta|^2
   =f(2-f)|\dd_\gamma\log s|^2\ge0,
 \qquad f=-s\beta_s.
\]
Differentiating $l^2+k_1^2=1$ gives $K=-e(\log k_1)$ and then
\eqref{pre:weighted-divergence}. The new normal and level second
fundamental form are $le$ and $lJ$, respectively. The new value of $S$ is
\[
 l(S+e(\beta))=lM_1+K/l.
\]
Substitute these expressions into the first three terms of
\eqref{pre:foliation-scalar}. Using $e(l)=\theta K/l$, their difference is
\[
 \begin{split}
 &-l^2|J|^2-(lM_1+K/l)^2-2l\,e(lM_1+K/l)\\
 &\hspace{14mm}+|J|^2+(M_1+K)^2+2e(M_1+K)\\
 &\qquad=\theta\bigl(|J|^2+M_1^2+K^2/l^2-2M_1K+2e(M_1)\bigr).
 \end{split}
\]
In particular, the derivatives of $K$ cancel. Adding the nonnegative
tangential change proves \eqref{pre:scalar-comparison}.
\end{proof}

Identity~\eqref{pre:weighted-divergence} explains the role of $M_1$:
it is the weighted divergence of $k_1e$, divided by $k_1$.
For the ordinary graph choice $\theta=s^2/(1+s^2)$, the vector field is
$k_1e=\nabla_hu/\sqrt{1+|\dd u|_h^2}$.
If this divergence is prescribed to be $p$, then $M_1=p/k_1$ and the
term $\theta M_1^2$ in \eqref{pre:scalar-comparison} is precisely $p^2$.
The terms involving $K$ and $e(M_1)$ still have to be controlled.
For the graph just constructed, the linear discrepancy pressure will
control them through the following square completion.

\subsection{Completion of the reserve estimate}

We can now turn the discrepancy pressure into the square in
Proposition~\ref{pre:reserve}. The argument is pointwise; consequently the
compact-data constants needed to construct $u$ will not enter the
reserve constant.

\begin{proof}[Proof of Proposition~\ref{pre:reserve}]
Take $u$ from Lemma~\ref{pre:global-graph}, set $s=|\dd u|_h$, and use
Lemma~\ref{pre:graph-comparison} with $t=w$ and
\[
 \theta=\frac{s^2}{1+s^2},\qquad
 l=(1+s^2)^{-1/2},\qquad k_1=ls,\qquad\beta=\log l.
\]
Here $0\le-s\beta_s\le1$. Define
\begin{equation}\label{pre:reserve-output}
 h^+=h+\dd u\otimes\dd u=h_u,
 \qquad t^+=w-\beta-u=t_u-u.
\end{equation}
Both expressions are smooth also where $\dd u=0$.

Write $p=c(u-w)$, with $c=10$. The weighted volume densities satisfy
\[
 e^{-t_u}\dd\vol_{h_u}=e^{-w}\dd\vol_h.
\]
Since $\nabla_{h_u}u=lk_1e$, this gives, where $s>0$,
\begin{equation}\label{pre:graph-laplacian}
 \Delta_{h_u}u-\langle\dd t_u,\dd u\rangle_{h_u}
   =\diver_{h,w}(lk_1e)
   =lp+\frac{k_1\theta K}{l},
 \qquad |\dd u|_{h_u}^2=k_1^2.
\end{equation}
Equation \eqref{pre:global-equation} and
\eqref{pre:weighted-divergence} give $M_1=p/k_1$, and hence
\begin{equation}\label{pre:pressure-derivative}
 e(M_1)=M_1K+\frac{c(s-e(w))}{k_1},
 \qquad e(w)=H-M_1-K.
\end{equation}
Subtracting $u$ from $t_u$ changes $\Dcal$ by
$-2\Delta_{h_u}u+2\langle\dd t_u,\dd u\rangle_{h_u}
-|\dd u|_{h_u}^2$. Substituting
\eqref{pre:graph-laplacian}--\eqref{pre:pressure-derivative} into
\eqref{pre:scalar-comparison} therefore yields
\begin{equation}\label{pre:reserve-expansion}
 \begin{split}
 \Dcal(h^+,t^+)-\Dcal(h,w)\ge{}&
 k_1^2|J|^2+p^2+\frac{k_1^2K^2}{l^2}
       +2ck_1(s-H+M_1+K)\\
 &-2\left(lp+\frac{k_1\theta K}{l}\right)-k_1^2.
 \end{split}
\end{equation}
There are three square completions. Since the level sets have dimension
$n-1$,
\[
 \begin{split}
 k_1^2|J|^2-2ck_1H&\ge-(n-1)c^2,\\
 \frac{k_1^2K^2}{l^2}+2k_1(c-\theta/l)K&\ge-(cl-\theta)^2,\\
 p^2+2(c-l)p&\ge\tfrac12p^2-2(c-l)^2.
 \end{split}
\]
For $c=10$, $0<l\le1$, and $0\le\theta<1$, we have
$|cl-\theta|\le10$, $|c-l|\le10$, and $k_1^2\le1$.
Thus \eqref{pre:reserve-expansion} implies
\begin{equation}\label{pre:reserve-intermediate}
 \Dcal(h^+,t^+)-\Dcal(h,w)
       \ge\tfrac12p^2+20k_1s-(100n+201).
\end{equation}
It remains to compare the right-hand side with $(t^+)^2$.
Since $t^+=-p/c-\beta$, the elementary inequality
\[
 -\beta=\tfrac12\log(1+s^2)\le\log(1+s)\le\sqrt s
\]
gives
\[
 (t^+)^2\le\frac{2p^2}{c^2}+2s.
\]
For $s\ge1$, $k_1\ge1/\sqrt2$, so $20k_1s\ge2s$.
For $s<1$, the term $2s$ costs at most $2$.
Together with \eqref{pre:reserve-intermediate}, this proves
\eqref{pre:reserve-bound} with $C(n)=100n+203$ wherever $s>0$.

The resulting inequality extends to boundary points of $\{s>0\}$
by smoothness. On the interior of $\{s=0\}$, the function $u$ is
locally constant and \eqref{pre:global-equation} gives $u=w$ there;
thus $t^+=0$ and the inequality holds directly. Finally, $h^+\ge h$
implies completeness: an $h^+$-Cauchy sequence is $h$-Cauchy, and its
$h$-limit is also its $h^+$-limit by local equivalence of smooth metrics.
\end{proof}

The new function $t^+$ may have either sign and need not be bounded
below. Proposition~\ref{pre:reserve} controls its square with a constant
depending only on the dimension. The local enlargement must therefore
use a pressure law that works for every real value of the drift
parameter; this is the purpose of the next section.
